\documentclass[10pt,reqno]{amsart}
\usepackage[T1]{fontenc}
\usepackage{lmodern,microtype,amsmath,amssymb,mathtools,amsthm}
\usepackage{enumitem,booktabs,longtable,array,geometry,hyperref}
\usepackage[nameinlink,capitalise,noabbrev]{cleveref}
\geometry{margin=0.91in}
\setlist{itemsep=1pt,topsep=3pt}
\setlength{\parskip}{0pt}
\emergencystretch=2em
\allowdisplaybreaks
\hypersetup{hidelinks,pdfauthor={cyclemath (editor and public releaser)},pdftitle={Convex Hulls of Great Subspheres and Gaussian Minima}}
\newtheorem{theorem}{Theorem}[section]
\newtheorem{proposition}[theorem]{Proposition}
\newtheorem{lemma}[theorem]{Lemma}
\newtheorem{corollary}[theorem]{Corollary}
\theoremstyle{definition}
\newtheorem{definition}[theorem]{Definition}
\theoremstyle{remark}
\newtheorem{remark}[theorem]{Remark}
\newcolumntype{P}[1]{>{\raggedright\arraybackslash}p{#1}}
\newcommand{\R}{\mathbb R}
\newcommand{\E}{\mathbb E}
\newcommand{\Prob}{\mathbb P}
\newcommand{\Sphere}{\mathbb S}
\newcommand{\conv}{\operatorname{conv}}
\newcommand{\Vol}{\operatorname{Vol}}
\newcommand{\rank}{\operatorname{rank}}
\newcommand{\Span}{\operatorname{span}}
\newcommand{\aff}{\operatorname{aff}}
\newcommand{\dist}{\operatorname{dist}}
\newcommand{\MA}{\operatorname{MA}}
\newcommand{\dd}{\,\mathrm d}
\newcommand{\ip}[2]{\langle #1,#2\rangle}
\title[Great Subspheres and Gaussian Minima]{Convex Hulls of Great Subspheres and Gaussian Minima}
\author{Edited and released by \MakeLowercase{cyclemath}}
\date{September 16, 2026}
\keywords{great subsphere, convex hull, Gaussian minimum, intrinsic volume, active stratum, sine polarity}
\begin{document}
\begin{abstract}
For unit normals $\nu_1,\ldots,\nu_n\in\Sphere^{d-1}$, $d\ge3$, let $K$ be the convex hull of the great subspheres $\Sphere^{d-1}\cap\nu_i^\perp$, and let $M=\min_i|\ip{G}{\nu_i}|$ for $G\sim N(0,I_d)$. The manuscript proves
\[
\Vol_d(K)=\kappa_d(1-\E M^2)-R_{n,d},\qquad R_{n,d}\ge0.
\]
In signed affine general position, the correction is a sum of nonnegative active-stratum integrals. The manuscript justifies its extension to every configuration and proves $R_{n,d}=0$ when the normals span at most a plane. In dimension three it also gives a direct formula for arbitrary configurations, including multiple ties: each antipodal pair of positive active polygons of area $A$ contributes $2A(1-r)^2/(3r)$, where $r$ is their circumcircle radius. For a triangle, writing $L=2A$, this is $L(1-r)^2/(3r)$. It establishes $V_2(K)=\pi(d-1-\E M^2)$ and deduces a conditional sharp volume bound from the $p=2$ part of Kunisky's Gaussian-minimum conjecture.
\end{abstract}
\maketitle
\noindent\textit{Status and attribution.} This is an AI-assisted research manuscript. Generative AI, principally OpenAI's ChatGPT, was used to carry out a substantial majority of the mathematical exploration, conjecture formation, proof development, proof revision, and diagnostic checking underlying this manuscript. cyclemath formulated and directed the investigation, selected and challenged intermediate claims, requested repeated audits and corrections, and prepared the manuscript for public release. The mathematical arguments have not received independent verification by a subject-matter expert, and the manuscript has not undergone peer review. Numerical checks in the supplementary material are diagnostics, not proofs.

\section{The statements and their logical scope}\label{sec:introduction}
Throughout the main theorem, $d\ge3$, $1\le n<\infty$, and every $\nu_i$ is a Euclidean unit vector. Repetitions and opposite normals are allowed. Write $B^d$ for the unit ball, $\kappa_d=\Vol_d(B^d)$, and $\dd\sigma$ for unnormalized surface measure. Thus $|\Sphere^{d-1}|=d\kappa_d$. Set
\begin{equation}\label{eq:disks}
D_i=B^d\cap\nu_i^\perp,\qquad
K=\conv\Bigl(\bigcup_{i=1}^nD_i\Bigr)
=\conv\Bigl(\bigcup_{i=1}^n(\Sphere^{d-1}\cap\nu_i^\perp)\Bigr).
\end{equation}
The last equality holds because each central unit disk is the convex hull of its boundary; it also holds for $d=2$, where the boundary is a pair of points. Let
\[
\Sigma=(\ip{\nu_i}{\nu_j})_{i,j},\quad
M(x)=\min_i|\ip{x}{\nu_i}|,\quad M=M(G),\quad G\sim N(0,I_d).
\]
Then $M$ has the Gaussian-minimum law associated with the correlation matrix $\Sigma$, and $\rank\Sigma=\dim\Span\{\nu_i\}$. Conversely every correlation matrix of rank at most $d$ is such a Gram matrix, by its positive semidefinite factorization.

\begin{theorem}[Volume--Gaussian identity]\label{thm:main}
For every configuration specified above,
\begin{equation}\label{eq:main-identity}
\boxed{\Vol_d(K)=\kappa_d(1-\E M^2)-R_{n,d},\qquad R_{n,d}\ge0.}
\end{equation}
The correction is the continuous function
$R_{n,d}=\kappa_d(1-\E M^2)-\Vol_d(K)$.
Under the sufficient general-position condition in \cref{def:agp}, it has the explicit representation \eqref{eq:R-explicit}. For every configuration whose normals span a subspace of dimension at most two, $R_{n,d}=0$. In dimension three, \eqref{eq:polygon-R} gives an explicit formula without general position.
\end{theorem}

\begin{proposition}[Second intrinsic volume]\label{prop:V2}
For $d\ge2$, $n\ge1$, and all unit-normal configurations,
\begin{equation}\label{eq:V2-gaussian}
\boxed{V_2(K)=\pi(d-1-\E M^2).}
\end{equation}
Here intrinsic volumes have Steiner normalization:
$\Vol_d(L+tB^d)=\sum_{j=0}^d\kappa_{d-j}V_j(L)t^{d-j}$, with $\kappa_0=1$.
In particular $V_2(B^d)=\pi(d-1)$, and the inequality in \cref{thm:main} is equivalent, for these bodies, to
\begin{equation}\label{eq:Vd-V2}
\Vol_d(K)\le\frac{\kappa_d}{\pi}\bigl(V_2(K)-\pi(d-2)\bigr).
\end{equation}
\end{proposition}

\begin{corollary}[Conditional sharp volume bound]\label{cor:kunisky}
Fix $n\ge1$ and $d\ge3$, and put
\begin{equation}\label{eq:cn}
c_n=1-\frac n\pi\sin\frac\pi n.
\end{equation}
If $\E M(\Sigma)^2\ge c_n$ for every $n\times n$ correlation matrix of rank at most $d$, then
\begin{equation}\label{eq:sharp-volume-bound}
\boxed{\Vol_d(K)\le\kappa_d\frac n\pi\sin\frac\pi n.}
\end{equation}
The equally spaced planar normal lines attain the displayed value, unconditionally. For $n=1$ both sides are zero.
\end{corollary}
\begin{proof}
Apply \cref{thm:main} to obtain
$\Vol_d(K)\le\kappa_d(1-\E M^2)\le\kappa_d(1-c_n)$.
The equality example is computed independently in \cref{sec:cosine}.
\end{proof}
For a particular configuration the volume bound is \emph{equivalent} to
\begin{equation}\label{eq:corrected-gaussian}
\E M^2+R_{n,d}/\kappa_d\ge c_n.
\end{equation}
The uncorrected moment inequality is a \emph{sufficient} condition. No converse to that implication is asserted. Under the moment hypothesis, equality in the volume bound holds exactly when both $\E M^2=c_n$ and $R_{n,d}=0$. Rank at most two is a sufficient condition for $R_{n,d}=0$; no rank characterization or uniqueness theorem for maximizers is asserted here.

\subsection*{Main contributions}
The principal result is \cref{thm:main}: for convex hulls of finitely many great subspheres in arbitrary ambient dimension, the ordinary volume is bounded by the second moment of the associated Gaussian minimum, with a nonnegative correction that is explicit in signed affine general position and extends continuously to every configuration. The proof organizes this correction by active strata and transfers flux successively from $2$-active through $d$-active strata by Stokes-type identities. In dimension three, \cref{prop:polygon} gives a direct active-polygon formula for arbitrary configurations, including symmetric multiple ties; under general position it reduces to the signed-triple expression \eqref{eq:d3-correction}. The manuscript also establishes the exact intrinsic-volume identity \eqref{eq:V2-gaussian} and the reduction to the span of the normals in \cref{prop:rank-reduction}. Finally, \cref{cor:kunisky} gives the resulting conditional sharp volume bound under the $p=2$ part of Kunisky's Gaussian comparison.

\subsection*{Relation to prior work and scope of the contribution}
Several parts of the surrounding framework are established in the literature and are separated here from the contribution of the present paper.

First, Kunisky introduced the cosine covariance, formulated the Gaussian stochastic comparison used here, and proposed the volumetric zone conjecture that would imply it \cite{Kunisky}. In Kunisky's formulation, the geometric conjecture concerns the surface measure of unions of spherical zones, rather than the ordinary Euclidean volume of the convex hull considered here. Gaussian minima and order statistics have a broader earlier literature; see, for example, Gordon, Litvak, Sch\"utt, and Werner \cite{GordonMinima,GordonMinimum} and Litvak's survey \cite{LitvakSurvey}. Those works provide estimates and extremal conjectures of a different type; the use of the Gaussian minimum here is specifically through Kunisky's covariance optimization problem.

Second, the polar description by intersections of origin-symmetric circular cylinders lies within the sine-polar framework of Huang, Li, Xi, and Ye \cite{HuangLiXiYe}. In the notation used below, if $P=\conv\{\pm\nu_1,\dots,\pm\nu_n\}$ and $[x,y]=\sqrt{|x|^2|y|^2-\ip{x}{y}^2}$, then $h_K(x)=\max_{y\in P}[x,y]$ and $K^\circ=P^\diamond$. Thus neither sine polarity nor the finite intersection-of-cylinders representation is claimed as new here. Explicit volumes of particular symmetric intersections of circular cylinders also have a much older literature; see, for example, Moore \cite{MooreCylinders}.

Third, there is a substantial classical theory relating intrinsic volumes to inradii, circumradii, successive radii, and tangential bodies; see Schneider \cite{Schneider} and, for bounds in terms of successive inner and outer radii, Henk and Hern\'andez Cifre \cite{HenkHernandez}. Those results provide important background but have a different form from \eqref{eq:Vd-V2}. In particular, no linear $V_d$--$V_2$ inequality for arbitrary convex bodies is asserted. Classical convex duality and subgradient conventions are used in the standard form of Rockafellar \cite{Rockafellar}.

Against this background, the contribution claimed here is deliberately specific. For the bodies obtained as convex hulls of finitely many great subspheres, the manuscript establishes \eqref{eq:main-identity}, equivalently the volume comparison \eqref{eq:Vd-V2}, and identifies its nonnegative active-stratum correction. It further gives the rank reduction \cref{prop:rank-reduction} and, in dimension three, the nongeneric active-polygon formula \cref{prop:polygon}, which specializes under general position to \eqref{eq:d3-correction}. It then deduces the conditional sharp volume bound from Kunisky's $p=2$ Gaussian comparison. Based on the literature search performed for this manuscript, and after comparison with the cylinder/sine-polar, intrinsic-volume, successive-radii, spherical-zone, and Gaussian-minimum literature described above, these volume--Gaussian results have not previously appeared in this form. This is a statement of literature awareness, not a claim that the absence of an equivalent reformulation elsewhere can be formally certified or that the literature search is exhaustive.

\subsection*{Guide to the proof}
The proof has four main ingredients.
\begin{enumerate}[label=(\roman*)]
\item The support function satisfies $h_K(x)^2=|x|^2-M(x)^2$. This yields the exact formula \eqref{eq:V2-gaussian} for the second intrinsic volume, after a support-function argument valid also for nonsmooth bodies.
\item Under signed affine general position, the ordinary volume $\Vol_d(K)$ is written as a sum over positive active strata using the convex potential $\Phi=\tfrac12h_K^2$. Convex duality supplies coverage, and the area formula gives multiplicity one almost everywhere.
\item The quantity $\kappa_d(1-\E M^2)$ is written as a Gaussian flux through pairwise interfaces. A cellwise spherical divergence identity and the oriented simplex face-flux identity transfer the discrepancy from level $k$ to level $k+1$, leaving the nonnegative scalar remainder $q_{d,k}$.
\item The recursion telescopes to \eqref{eq:R-explicit}. Continuity removes the general-position assumption for the total correction. A separate rank-reduction argument proves the rank-at-most-two equality, and in $d=3$ a polygon version of the same flux mechanism treats multiple ties directly.
\end{enumerate}

\section{Support functions, dimension, and the intrinsic-volume identity}\label{sec:prelim}
For any compact convex set $L$, write $h_L(x)=\sup_{y\in L}\ip{x}{y}$. For the present bodies,
\begin{equation}\label{eq:support}
h_K(x)^2=\max_i|P_{\nu_i^\perp}x|^2=|x|^2-M(x)^2.
\end{equation}
In particular, with $m(u)=M(u)$ on $\Sphere^{d-1}$,
\begin{equation}\label{eq:h-m}
h_K(u)^2=1-m(u)^2,\qquad
K^\circ=\bigcap_i\{x:|P_{\nu_i^\perp}x|\le1\}.
\end{equation}
If $P=\conv\{\pm\nu_i\}$ and $[x,y]=\sqrt{|x|^2|y|^2-\ip{x}{y}^2}$, the function $y\mapsto[x,y]$ is a seminorm. Its maximum on $P$ is its maximum on the listed vertices, so
$h_K(x)=\max_{y\in P}[x,y]$ and $K^\circ=P^\diamond$, with the sine-polar convention of \cite{HuangLiXiYe}. These set identities remain valid for lower-dimensional $P$; boundedness of the polar is a separate issue.

\begin{lemma}[Dimension and repeated lines]\label{lem:dimension}
For $d\ge2$, $K$ has nonempty interior in $\R^d$ if and only if at least two normal lines are different, equivalently $\rank\Sigma\ge2$. If $\rank\Sigma=1$, then $K$ is a unit $(d-1)$-ball in a hyperplane, $\Vol_d(K)=0$, and $\E M^2=1$. Deleting repeated normal lines does not change $K$, $M(x)$, or $R_{n,d}$.
\end{lemma}
\begin{proof}
In rank one all disks coincide. Otherwise choose $\nu_i\ne\pm\nu_j$.
The two hyperplanes $\nu_i^\perp$ and $\nu_j^\perp$ span $\R^d$, so the symmetric convex hull of their unit balls spans $\R^d$ and contains a neighborhood of zero. The last assertion follows directly from the defining maximum and minimum. Deletion changes the number of entries $n$, so it does not preserve a conjectural comparison constant indexed by $n$; the original $n$ is retained in such comparisons.
\end{proof}

\begin{lemma}[The support-function quadratic formula]\label{lem:V2-support}
For every compact convex $L\subset\R^d$, $d\ge2$,
\begin{equation}\label{eq:V2-support}
V_2(L)=\frac{\pi}{d\kappa_d}\int_{\Sphere^{d-1}}
\bigl((d-1)h_L^2-|\nabla_{\Sphere}h_L|^2\bigr)\dd\sigma,
\end{equation}
where the gradient is the almost-everywhere, equivalently weak, spherical gradient.
\end{lemma}
\begin{proof}
First take a smooth support function with positive curvature. Kubota's formula in Steiner normalization \cite{Schneider} gives
$V_2(L)=(d-1)\E_F\Vol_2(P_FL)$ for a uniformly random two-plane $F$.
The planar support-area formula is
$\Vol_2(P_FL)=\frac12\int_{\Sphere(F)}(h_L^2-(\partial_t h_L)^2)\dd\theta$;
it follows by parametrizing the boundary as $hu+h't$ and integrating the signed area.
A random plane and a random direction $u$ in it give uniform $u\in\Sphere^{d-1}$ and uniform unit tangent $t\in u^\perp$.
Thus $\E_t(\partial_t h_L)^2=|\nabla_{\Sphere}h_L|^2/(d-1)$.
Substitution gives \eqref{eq:V2-support}, including its constant.

Here is the extension argument, including the gradient term. Average rotated copies of $L$ against a smooth approximate identity on $SO(d)$ and add a ball of radius tending to zero. This gives smooth, strictly convex $L_j$ with $h_j\to h_L$ uniformly and all $L_j$ contained in one fixed ball. Let $H_j,H$ be their homogeneous support functions on $\R^d$. At a direction $u$ where $H$ is differentiable, $\nabla H_j(u)$ is bounded. Every subsequential limit $p$ satisfies
\[
H(x)\ge H(u)+\ip{p}{x-u}\quad(x\in\R^d)
\]
by the corresponding subgradient inequality for $H_j$. Hence $p=\nabla H(u)$.
Convex Lipschitz functions are differentiable almost everywhere; homogeneity transfers this conclusion to almost every sphere direction. Therefore $\nabla_{\Sphere}h_j\to\nabla_{\Sphere}h_L$ almost everywhere, with a common bound. Dominated convergence yields convergence of the squared-gradient integrals. Intrinsic volumes converge under Hausdorff convergence: this also follows by evaluating the Steiner polynomials at $d+1$ fixed positive values of $t$ and passing to their coefficients. This proves the formula, even when $L$ has empty interior.
\end{proof}

\begin{proof}[Proof of \cref{prop:V2}]
First delete duplicate lines. For different remaining lines, the equality
$|\ip{u}{\nu_i}|=|\ip{u}{\nu_j}|$ is contained in the union of the two proper great hyperplanes orthogonal to $\nu_i-\nu_j$ and $\nu_i+\nu_j$. Thus a unique quadratic piece is active almost everywhere. Away also from the finite set where a selected projection has absolute value one, put $s=\ip{u}{\nu_i}$; then
\[
h_K=\sqrt{1-s^2},\quad
\nabla_{\Sphere}s=\nu_i-su,\quad
|\nabla_{\Sphere}h_K|^2=s^2=m(u)^2.
\]
Consequently $(d-1)h_K^2-|\nabla_{\Sphere}h_K|^2=d-1-dm(u)^2$ almost everywhere.
For $G=\rho U$, the variable $U$ is uniform on the sphere, independent of $\rho$, with $\E\rho^2=d$. Therefore
\begin{equation}\label{eq:radial-moment}
\E M^2=d\E m(U)^2=\frac1{\kappa_d}\int_{\Sphere^{d-1}}m(u)^2\dd\sigma(u).
\end{equation}
Insert this in \eqref{eq:V2-support}. Reinstating duplicate labels does not change either side.
\end{proof}
Unique \emph{index} attainment almost everywhere holds if and only if no line is repeated. For necessity, take a generic direction in the hyperplane perpendicular to a repeated line, avoiding the other finitely many hyperplanes; in a neighborhood the repeated line is the strict minimizing line. Uniqueness of the index is not required in \cref{prop:V2}.

\begin{remark}[Special structure of the volume inequality]\label{rem:special-structure}
The proof of \cref{prop:V2} uses only the support function. By contrast, \eqref{eq:Vd-V2} is not asserted as a general intrinsic-volume inequality. The proof of \cref{thm:main} uses the special representation \eqref{eq:disks}: every generator is a full central hyperplane section of the Euclidean ball.
\end{remark}

\section{Signed cells and their dimensions}\label{sec:active}
\begin{definition}[Signed affine general position]\label{def:agp}
The signed labels are $\mathcal S=\{(i,\epsilon):1\le i\le n,\ \epsilon\in\{-1,1\}\}$, with associated vectors $a_{i,\epsilon}=\epsilon\nu_i$. A signed set contains at most one label of each original index. The configuration is in signed affine general position (AGP) if every such set of size $k\le\min(n,d+1)$ has affinely independent associated vectors.
\end{definition}
AGP is a \emph{sufficient device} for the stratum proofs, not a necessary condition for the main theorem or for their numerical identities. It is open and dense in $(\Sphere^{d-1})^n$: each condition is positivity of the Gram determinant of $k-1$ difference vectors, a real-analytic function not identically zero, since $k\le d+1$ affinely independent points can be chosen on the sphere. Its zero set has empty interior on this connected analytic manifold. There are only finitely many conditions. For $n=1$ they are vacuous and the rank-one calculation in \cref{lem:dimension} applies.

For a signed set $\tau=\{a_1,\ldots,a_k\}$ of size $2\le k\le\min(n,d)$, the vector symbols are also used to denote their labels, without discarding the labels. Define
\begin{gather}
D_\tau=\Span\{a_2-a_1,\ldots,a_k-a_1\},\quad W_\tau=D_\tau^\perp,
\quad w_\tau=P_{W_\tau}a_1,\label{eq:D-W}\\
c_\tau=|w_\tau|,\quad \mu_\tau(\theta)=\ip{w_\tau}{\theta},\quad
\Delta_\tau=\conv\{a_1,\ldots,a_k\},\quad A_\tau=\Vol_{k-1}(\Delta_\tau),\label{eq:cell-params}\\
\Omega_\tau=\{\theta\in\Sphere(W_\tau):\mu_\tau(\theta)>0,
\ |\ip{\nu_\ell}{\theta}|>\mu_\tau(\theta)\text{ for all inactive }\ell\}.
\label{eq:active-cell}
\end{gather}
Write $\mathcal A_k=\{\tau:\Omega_\tau\ne\varnothing\}$. Sets $\tau$ and $-\tau$ are counted separately. All active vectors have projection $w_\tau$ onto $W_\tau$, and $w_\tau$ is the point of $\aff\Delta_\tau$ nearest zero.

\begin{lemma}[Geometry required by integration]\label{lem:cell-geometry}
Under AGP, each nonempty $\Omega_\tau$ has dimension $d-k$, $A_\tau>0$, and $0<c_\tau<1$. On its closure, $0\le\mu_\tau\le c_\tau$.
Within $\mu_\tau>0$, it is a finite union of spherical polyhedral cells with corners. Each boundary hypersurface adds exactly one signed label; simultaneous additions of $j$ labels have codimension $j$ when nonempty. No positive active set has more than $d$ labels. For $k=d$, $\Omega_\tau$ consists of the single positive point $w_\tau/c_\tau$ and its measure is counting measure.
\end{lemma}
\begin{proof}
Affine independence gives $\dim D_\tau=k-1$ and $\dim W_\tau=d-k+1$. The strict inequalities define an open subset of its sphere. Nonemptiness and $\mu>0$ imply $c>0$. Write $a_j=w+z_j$, $z_j\in D_\tau$. Then $|z_j|^2=1-c^2$; if $c=1$, all $a_j$ coincide, impossible for $k\ge2$.

On $\mu>0$ the signs of every inactive projection are fixed in each component. After choosing them, all inequalities are linear inequalities of the form $\ip{a_\ell-a_1}{\theta}>0$. Adding $j$ labels adds $j$ independent differences modulo $D_\tau$, by AGP. At their common zero set, those projected differences are tangent to $\Sphere(W_\tau)$, so the boundary gradients are independent. If $d+1$ labels were active, $d$ independent differences would annihilate the nonzero vector $\theta$, impossible. These observations prove the dimension and boundary assertions without a transversality assumption left implicit.
\end{proof}

\section{A volume formula with coverage and multiplicity}\label{sec:volume}
\begin{lemma}[Conjugacy and level matching]\label{lem:duality}
Suppose $K$ has nonempty interior and is symmetric about zero. Put $\Phi=h_K^2/2$ and let $p_K$ be its gauge. Then
\begin{equation}\label{eq:duality}
\Phi^*(y)=p_K(y)^2/2,\qquad
y\in\partial\Phi(x)\ \Longrightarrow\ p_K(y)=h_K(x),\qquad
\partial\Phi(K^\circ)=K.
\end{equation}
For almost every $y$, at most one $x$ satisfies $y\in\partial\Phi(x)$.
\end{lemma}
\begin{proof}
Write $x=t z$ with $h_K(z)=1$, $t\ge0$. Maximizing $\ip{x}{y}-h_K(x)^2/2$ first over $z\in K^\circ$ and then over $t$ gives
$\sup_{t\ge0}(tp_K(y)-t^2/2)=p_K(y)^2/2$.
Both sides are finite because zero is an interior point of $K$. If $y\in\partial\Phi(x)$, the one-dimensional subgradient inequality on $tx$ gives $\ip{x}{y}=2\Phi(x)$. Fenchel equality gives $\Phi^*(y)=\Phi(x)$ and hence level matching.
For each $y$, the conjugate maximum is attained, since $h_K(x)\ge a|x|$ for some $a>0$ makes the maximized function coercive downwards. The maximizer $x$ satisfies $y\in\partial\Phi(x)$; level matching shows it belongs to $K^\circ$ exactly when $y\in K$.
Finally, $y\in\partial\Phi(x)$ is equivalent to $x\in\partial\Phi^*(y)$. The finite convex function $\Phi^*$ is locally Lipschitz and differentiable almost everywhere; at those $y$, $x=\nabla\Phi^*(y)$ is unique.
\end{proof}
In rank one, the correct gauge is extended-valued: if $K=B^d\cap\nu^\perp$, then $p_K(y)=|y|$ on $\nu^\perp$ and $+\infty$ elsewhere. Direct conjugation gives the same extended-valued identity. The full-dimensional proof is not applied to that case.

\begin{proposition}[Active-stratum volume]\label{prop:active-volume}
For AGP configurations with $n\ge2$,
\begin{equation}\label{eq:active-volume}
\boxed{\Vol_d(K)=\frac1d\sum_{k=2}^{\min(n,d)}\sum_{\tau\in\mathcal A_k}
A_\tau(1-c_\tau^2)\int_{\Omega_\tau}
\frac{\mu_\tau^{k-1}}{(1-\mu_\tau^2)^{d/2}}\dd\sigma.}
\end{equation}
\end{proposition}
\begin{proof}
The convex potential is a finite maximum of quadratics:
\[
\Phi(x)=\tfrac12\max_i(|x|^2-\ip{x}{\nu_i}^2).
\]
For $x=\rho\theta$, $\rho>0$, $\theta\in\Omega_\tau$, the subdifferential of a finite maximum is the convex hull of the active gradients; equivalently this follows by taking directional derivatives of that maximum. Thus
\begin{equation}\label{eq:subgradient-simplex}
\partial\Phi(x)=\rho(\theta-\mu_\tau\Delta_\tau).
\end{equation}
Let $z\in\Delta_\tau-w_\tau\subset D_\tau$, and write $w=w_\tau$. Relative to $W_\tau\oplus D_\tau$, the image point $y$ is
\begin{equation}\label{eq:block-map}
y_W=(I-w\otimes w)x,\qquad y_D=-\ip{w}{x}z.
\end{equation}
Since $c<1$ and $\ip{w}{x}>0$, this parametrization is injective for the fixed stratum: recover $x$ from $y_W$, then $z$ from $y_D$. In the product measures $\dd x\,\dd z$, its absolute Jacobian is
$(1-c^2)\ip{w}{x}^{k-1}$. Polar coordinates in $W_\tau$, whose dimension is $d-k+1$, give
\[
(1-c^2)\mu^{k-1}\rho^{d-1}\dd\rho\,\dd\sigma\,\dd z.
\]
The condition $x\in K^\circ$ is $0\le\rho\le(1-\mu^2)^{-1/2}$. Integrating in $z$ gives $A_\tau$, and integrating in $\rho$ gives the summand in \eqref{eq:active-volume}.

It remains to justify summation as volume, rather than volume with multiplicity. By \cref{lem:duality}, all of $K$ is covered by these subgradient images together with the following omitted pieces. At $x=0$ the image is $0$. If $M(x)=0$, every active gradient equals $x$, so the image lies in the finite union $\bigcup_i\nu_i^\perp$. If exactly one index is active, its gradient is $P_{\nu_i^\perp}x$, again in a fixed hyperplane. These images all have $d$-dimensional measure zero. Every other $x$ has exactly one of the positive signed active sets in \cref{lem:cell-geometry}. Finally, different $x$ can map to the same $y$ only in the null set where $\Phi^*$ is not differentiable. For the same $x$, the active set and the parameter $z$ are unique. The area formula, applied to these finitely many bounded Lipschitz parametrizations (or compact exhaustions), therefore sums their volumes with multiplicity one almost everywhere. This proves the claimed equality. It also justifies the Alexandrov notation $\MA_\Phi(K^\circ)=|\partial\Phi(K^\circ)|=\Vol_d(K)$, if desired.
\end{proof}

\section{Gaussian pair flux and the Stokes recursion}\label{sec:stokes}
\begin{proposition}[Gaussian pair flux]\label{prop:gaussian-flux}
For an AGP configuration with $n\ge2$,
\begin{equation}\label{eq:gaussian-flux}
\boxed{\kappa_d(1-\E M^2)=\frac1d\sum_{\tau\in\mathcal A_2}
A_\tau\int_{\Omega_\tau}\mu_\tau\dd\sigma.}
\end{equation}
\end{proposition}
\begin{proof}
For a full sign choice $s\in\{-1,1\}^n$, put $a_j=s_j\nu_j$. Subdivide space into the polyhedral cones
\[
T_{s,i}=\{x:0\le\ip{a_i}{x}\le\ip{a_j}{x}\text{ for all }j\}.
\]
Their full-dimensional interiors partition space away from finitely many hyperplanes; cones with empty interior have zero volume. With $\gamma_d(x)=(2\pi)^{-d/2}e^{-|x|^2/2}$, use on $T_{s,i}$ the vector field
\[
X_i(x)=\ip{a_i}{x}a_i\gamma_d(x),\qquad
\operatorname{div}X_i=(1-\ip{a_i}{x}^2)\gamma_d(x).
\]
Apply the divergence theorem on $T_{s,i}\cap B_R$. The artificial spherical-boundary flux is at most a constant times $R^d e^{-R^2/2}$ and tends to zero. Dominated convergence applies to the volume integrals. All sign-zero boundaries have $\ip{a_i}{x}=0$ and hence zero flux; higher active intersections have codimension at least two under AGP and no interface surface measure.

On a positive interface between indices $i,j$, set $m=\ip{a_i}{x}=\ip{a_j}{x}>0$. The outward normal from $T_{s,i}$ is $(a_i-a_j)/|a_i-a_j|$ because its defining inequality is $\ip{a_i-a_j}{x}\le0$. The two outward fluxes sum to $m|a_i-a_j|\gamma_d(x)$. Inactive signs are uniquely determined there. Thus the interface is counted once as an unordered signed pair $\tau$, with $A_\tau=|a_i-a_j|$. Its cone has dimension $d-1$, so its radial integration is
\[
(2\pi)^{-d/2}\int_0^\infty\rho^{d-1}e^{-\rho^2/2}\dd\rho
=\frac1{d\kappa_d}.
\]
The additional power of $\rho$ comes from $m=\rho\mu_\tau$. Summing the remaining angular integrals proves \eqref{eq:gaussian-flux}.
\end{proof}

For $0\le t<1$, set $f_{d,1}(t)=t$ and define
\begin{equation}\label{eq:fdk}
f_{d,k}(t)=\int_0^t\frac{s^{k-1}}{(1-s^2)^{d/2}}\dd s\qquad(2\le k\le d-1).
\end{equation}
For $2\le k<d$, and separately for the terminal index, put
\begin{align}
q_{d,k}(t)&=(k-1)f_{d,k-1}(t)+(d-k)tf_{d,k}(t)
-t^{k-1}(1-t^2)^{1-d/2},\label{eq:qdk}\\
q_{d,d}(t)&=(d-1)f_{d,d-1}(t)-t^{d-1}(1-t^2)^{1-d/2}.\label{eq:qdd}
\end{align}
No $f_{d,d}$ is defined or needed. Empty sums below are zero.

\begin{lemma}[Scalar remainders]\label{lem:q-positive}
For every integer $d\ge3$, $q_{d,2}=0$ on $[0,1)$. For $3\le k\le d$, $q_{d,k}(0)=0$ and $q_{d,k}(t)>0$ for $0<t<1$.
\end{lemma}
\begin{proof}
Integration gives
\begin{equation}\label{eq:fd2}
f_{d,2}(t)=\frac{(1-t^2)^{1-d/2}-1}{d-2},
\end{equation}
which gives $q_{d,2}=0$. For $3\le k<d$, differentiation yields
\begin{equation}\label{eq:q-derivative}
q_{d,k}'(t)=(d-k)f_{d,k}(t)+t^k(1-t^2)^{-d/2}.
\end{equation}
Indeed, the remaining coefficients after extracting $t^k(1-t^2)^{-d/2}$ are
$(k-1)+(d-k)-(d-2)=1$. For $k=d$ the derivative of \eqref{eq:qdd} is just $t^d(1-t^2)^{-d/2}$. The derivative formula is not applied to $k=2$, where $f_{d,1}'=1$ is different from the rule for $k\ge2$. Positivity now follows by integration from zero.
\end{proof}

\begin{lemma}[Cell divergence and the zero boundary]\label{lem:cell-divergence}
Let $\tau\in\mathcal A_k$ in AGP, $2\le k\le d-1$. On $\Sphere(W_\tau)$, abbreviate $c=c_\tau$, $\mu=\mu_\tau$. Then
\begin{align}
|\nabla_{\Sphere}\mu|^2&=c^2-\mu^2,&
\Delta_{\Sphere}\mu&=-(d-k)\mu,\label{eq:spherical-linear}\\
(k-1)f_{d,k-1}(\mu)-(1-c^2)\frac{\mu^{k-1}}{(1-\mu^2)^{d/2}}
&=\operatorname{div}_{\Sphere(W_\tau)}(f_{d,k}(\mu)\nabla_{\Sphere}\mu)+q_{d,k}(\mu).
\label{eq:cell-divergence}
\end{align}
Integration of the divergence leaves precisely its positive active boundary fluxes: the boundary at $\mu=0$ contributes zero in the limit.
\end{lemma}
\begin{proof}
A linear function on a sphere of dimension $d-k$ has the gradient and Laplacian in \eqref{eq:spherical-linear}. The product rule gives
\[
\operatorname{div}(f_{d,k}(\mu)\nabla\mu)
=\frac{\mu^{k-1}(c^2-\mu^2)}{(1-\mu^2)^{d/2}}-(d-k)\mu f_{d,k}(\mu).
\]
Subtract this from the left side of \eqref{eq:cell-divergence} to get \eqref{eq:qdk}.
For the integration, cut each spherical polyhedral component by $\mu>\varepsilon$, choosing a sequence of regular values for its finitely many faces. Apply Stokes on this compact piecewise smooth region. All integrands are bounded for the fixed configuration because $\mu\le c<1$. On the artificial boundary,
\[
f_{d,k}(\varepsilon)=O(\varepsilon^k),\qquad
|\partial_n\mu|\le c.
\]
The total $(d-k-1)$-measure of a level section $\mu=\varepsilon$ stays bounded as $\varepsilon\downarrow0$: it is bounded by the measure of the full latitude sphere, times the finite number of components. If $d-k=1$ this is a bound on the number of endpoints. Hence the artificial flux is $O(\varepsilon^k)\to0$. The same bounds and the finite great-sphere faces justify convergence of the other integrals. This treats the zero boundary at every level, without assuming it is a regular face before truncation.
\end{proof}

\begin{lemma}[Oriented simplex face flux]\label{lem:simplex-flux}
Let $\eta=\{a_0,\ldots,a_k\}\in\mathcal A_{k+1}$ in AGP, $2\le k\le d-1$, and let $\tau_j=\eta\setminus\{a_j\}$. Every $\tau_j$ has a locally adjacent active cell along $\Omega_\eta$. With $n_j$ outward from that cell,
\begin{equation}\label{eq:simplex-flux}
\sum_{j=0}^k A_{\tau_j}\partial_{n_j}\mu_{\tau_j}=kA_\eta.
\end{equation}
No positivity assumption on barycentric coordinates is needed.
\end{lemma}
\begin{proof}
Let $w=w_\eta=\sum_{j=0}^k\lambda_j a_j$, $\sum_j\lambda_j=1$.
Choose the unit vector $e_j\in D_\eta\cap D_{\tau_j}^{\perp}$ pointing from the opposite face toward $a_j$, and let $h_j>0$ be the altitude. For each $i\ne j$,
$\ip{a_j-a_i}{e_j}=h_j$; since $w\perp D_\eta$, barycentric averaging gives
$\ip{a_i}{e_j}=-\lambda_jh_j$ for $i\ne j$. Consequently
\[
w_{\tau_j}=w-\lambda_jh_je_j.
\]
At $\theta\in\Omega_\eta$, $e_j\perp\theta$. Along
$\theta(t)=(\theta+te_j)/|\theta+te_j|$, the difference between the dropped value $\ip{a_j}{\theta(t)}$ and the common face value has derivative $h_j>0$ at zero. Positivity of the common value and all other strict gaps persist for small $t>0$. Thus this side is the interior of the face cell and $n_j=-e_j$. It follows that
$\partial_{n_j}\mu_{\tau_j}=\ip{w_{\tau_j}}{-e_j}=\lambda_jh_j$.
Using $A_{\tau_j}h_j=kA_\eta$ and summing proves \eqref{eq:simplex-flux}. Individual terms can be negative when $w$ lies outside the simplex; only their sum has the asserted positive sign.
\end{proof}

\begin{remark}[The recursive picture]\label{rem:recursive-picture}
The mechanism may be summarized schematically as
\[
\text{Gaussian pair flux}\longrightarrow \mathcal A_2\longrightarrow \mathcal A_3\longrightarrow\cdots\longrightarrow \mathcal A_d.
\]
At level $k$, the divergence term moves to level $k+1$, while $q_{d,k}\ge0$ remains behind as part of the correction. The oriented face-flux identity is what makes the coefficient at the next level exactly right.
\end{remark}

\begin{proof}[Proof of \cref{thm:main} under AGP]
The $n=1$ case was settled in \cref{lem:dimension}. Otherwise let $V_k$ denote the level-$k$ contribution to \eqref{eq:active-volume}, and put
\[
I_k=\frac1d\sum_{\tau\in\mathcal A_k}A_\tau
\int_{\Omega_\tau}(k-1)f_{d,k-1}(\mu_\tau)\dd\sigma,
\qquad
Q_k=\frac1d\sum_{\tau\in\mathcal A_k}A_\tau
\int_{\Omega_\tau}q_{d,k}(\mu_\tau)\dd\sigma.
\]
Indices $k>n$ give empty sums. By \cref{prop:gaussian-flux}, $I_2=\kappa_d(1-\E M^2)$.
For $2\le k<d$, \cref{lem:cell-divergence,lem:simplex-flux} give the exact recurrence
\begin{equation}\label{eq:recurrence}
I_k-V_k=Q_k+I_{k+1}.
\end{equation}
Indeed, each positive codimension-one face is a $(k+1)$-active stratum by \cref{lem:cell-geometry}, and the sum of its incoming coefficients is $kA_\eta$. Zero boundaries vanish and corners have no boundary measure. Since $Q_2=0$, telescoping \eqref{eq:recurrence} from $2$ to $d-1$ leaves $\sum_{k=3}^{d-1}Q_k+I_d-V_d$.
At a terminal stratum $\mu_\tau=c_\tau$, and its sphere is zero-dimensional; hence
\[
I_d-V_d=\frac1d\sum_{\tau\in\mathcal A_d}A_\tau q_{d,d}(c_\tau).
\]
It follows that
\begin{equation}\label{eq:R-explicit}
\boxed{R_{n,d}=\frac1d\sum_{k=3}^{d-1}\sum_{\tau\in\mathcal A_k}
A_\tau\int_{\Omega_\tau}q_{d,k}(\mu_\tau)\dd\sigma
+\frac1d\sum_{\tau\in\mathcal A_d}A_\tau q_{d,d}(c_\tau)\ge0.}
\end{equation}
Each term is nonnegative by \cref{lem:q-positive}. This proves the generic identity and inequality. It also shows, under AGP, that $R_{n,d}=0$ if and only if every $\mathcal A_k$, $3\le k\le d$, is empty: each nonempty open stratum has positive measure (counting measure at the terminal level) and strictly positive integrand and simplex volume.
\end{proof}

\section{Continuity, rank, and ambient dimension}\label{sec:degenerate}
\begin{proposition}[Continuous extension]\label{prop:continuity}
For fixed $n,d$, $R_{n,d}=\kappa_d(1-\E M^2)-\Vol_d(K)$ is continuous on $(\Sphere^{d-1})^n$. In particular, \eqref{eq:main-identity} holds without AGP.
\end{proposition}
\begin{proof}
A perturbation of the normals makes the support functions in \eqref{eq:support} converge uniformly on the sphere, hence makes $K$ converge in the Hausdorff metric. The $d$-volume converges, including when the limit is lower-dimensional. For the Gaussian part, $M_j(G)^2\to M(G)^2$ pointwise and $0\le M_j(G)^2\le|G|^2$, whose expectation is $d$. Dominated convergence proves continuity. Approximate by the dense AGP configurations and pass the nonnegative global difference to the limit. Individual stratum integrals or their labels need not have separate limits; no termwise limiting claim is used.
\end{proof}

\begin{proposition}[Reduction to the span of the normals]\label{prop:rank-reduction}
Let $L=\Span\{\nu_i\}$ have dimension $r\ge2$, let $E=L^\perp$, and define
$K_L=\conv\bigl(\bigcup_i(B^r_L\cap\nu_i^\perp)\bigr)\subset L$.
Then
\begin{gather}
h_K(x_L+x_E)^2=h_{K_L}(x_L)^2+|x_E|^2,\label{eq:rank-support}\\
K=\{y_L+y_E:p_{K_L}(y_L)^2+|y_E|^2\le1\},\qquad
\frac{\Vol_d(K)}{\kappa_d}=\frac{\Vol_r(K_L)}{\kappa_r}.
\label{eq:rank-volume}
\end{gather}
The Gaussian minimum has the same law in ambient dimensions $r$ and $d$. Consequently, for $r\ge3$, $R_{n,d}/\kappa_d=R_{n,r}/\kappa_r$.
\end{proposition}
\begin{proof}
The first identity follows by taking the maximum in
$|P_{\nu_i^\perp}(x_L+x_E)|^2=|x_E|^2+|x_L|^2-\ip{x_L}{\nu_i}^2$.
The body on the right of \eqref{eq:rank-volume} has support
$\sqrt{h_{K_L}(x_L)^2+|x_E|^2}$: maximize $a h_{K_L}(x_L)+b|x_E|$ over $a,b\ge0$, $a^2+b^2\le1$. Therefore it is $K$.
Its section at $y_E=z$, $|z|\le1$, is $\sqrt{1-|z|^2}K_L$. Fubini gives
\[
\Vol_d(K)=\Vol_r(K_L)\int_{B^{d-r}}(1-|z|^2)^{r/2}\dd z.
\]
Applying the same section calculation to the unit ball identifies the integral as $\kappa_d/\kappa_r$. This includes $r=d$, with integration over a zero-dimensional space interpreted as one. Finally $P_LG\sim N(0,I_L)$, and only this projection enters $M$. The statement about $R$ follows by subtraction.
\end{proof}

\begin{proposition}[Rank at most two]\label{prop:rank-two}
If $\rank\Sigma\le2$, then $R_{n,d}=0$.
\end{proposition}
\begin{proof}
Rank one is \cref{lem:dimension}. In rank two let $J$ be rotation by $\pi/2$ in $L$ and put $Q=\conv\{\pm J\nu_i\}\subset L$. Then $K_L=Q$.
The already proved $d=2$ version of \cref{prop:V2}, where $V_2(Q)=\Vol_2(Q)$, gives $\Vol_2(Q)=\pi(1-\E M^2)$. By \cref{prop:rank-reduction},
$\Vol_d(K)=(\kappa_d/\pi)\Vol_2(Q)=\kappa_d(1-\E M^2)$.
This proof uses no perturbation confined to a plane and no generic active-stratum formula on a nongeneric configuration.
\end{proof}
Together with \cref{prop:continuity}, this completes all parts of \cref{thm:main} except the separately proved three-dimensional formula below.

\section{Three dimensions, including symmetric multiple ties}\label{sec:d3}
For an AGP signed active triple $\tau=\{a,b,c\}$, put
\begin{equation}\label{eq:d3-params}
L_\tau=|(b-a)\times(c-a)|,\quad m_\tau=\dist(0,\aff\{a,b,c\}),\quad
r_\tau=\sqrt{1-m_\tau^2}.
\end{equation}
Here $A_\tau=L_\tau/2$, $0<m_\tau<1$, and $r_\tau>0$.
Since $f_{3,2}(m)=1/r-1$,
\begin{equation}\label{eq:d3-q}
q_{3,3}(m)=2(1/r-1)-m^2/r=(1-r)^2/r.
\end{equation}
One signed triple contributes $L_\tau(1-r_\tau)^2/(6r_\tau)$.
Thus, counting each antipodal pair exactly once,
\begin{equation}\label{eq:d3-correction}
R_{n,3}=\sum_{\{\tau,-\tau\}}\frac{L_\tau(1-r_\tau)^2}{3r_\tau}
\quad\text{under AGP}.
\end{equation}
The factor of two distinguishes a single signed triple from its antipodal pair.

\begin{proposition}[Active-polygon formula without general position]\label{prop:polygon}
Let $d=3$, delete repeated normal lines, and let $\mathcal V$ be the set of directions $\theta\in\Sphere^2$ for which $m(\theta)>0$ and at least three remaining indices attain the minimum. For $\theta\in\mathcal V$ define
\begin{gather*}
P_\theta=\conv\{\operatorname{sign}(\ip{\theta}{\nu_i})\nu_i:
|\ip{\theta}{\nu_i}|=m(\theta)\},\quad
A_\theta=\Vol_2(P_\theta),\quad r_\theta=\sqrt{1-m(\theta)^2}.
\end{gather*}
The set $\mathcal V$ is finite, every $P_\theta$ is a nondegenerate polygon, and
\begin{equation}\label{eq:polygon-R}
\boxed{R_{n,3}=\frac13\sum_{\theta\in\mathcal V}A_\theta\frac{(1-r_\theta)^2}{r_\theta}
=\frac23\sum_{\{\theta,-\theta\}\subset\mathcal V}A_\theta\frac{(1-r_\theta)^2}{r_\theta}.}
\end{equation}
Every antipodal pair in the last sum is counted once. The formula includes empty $\mathcal V$, in particular rank at most two. It implies $R_{n,3}=0$ if and only if $\mathcal V$ is empty.
\end{proposition}
\begin{proof}
Distinct points on a sphere cannot contain three collinear points. Three positive active signed vectors therefore determine a unique affine plane with a one-dimensional orthogonal space; the positive normal direction is unique. There are finitely many signed triples, proving finiteness of $\mathcal V$. All active vectors at $\theta$ lie on the circle with center $w=m(\theta)\theta$ and radius $r_\theta>0$ in the plane $\ip{\theta}{a}=m(\theta)$. Distinct points on this circle are vertices of a nondegenerate convex polygon.

If only one normal line remains, both sides of \eqref{eq:polygon-R} vanish. Otherwise $K$ has interior. The coverage and multiplicity proof in \cref{prop:active-volume} still applies: away from $M(x)=0$ and the unique-index regions, the positive strata are open pair arcs and the isolated directions $\mathcal V$. At a direction $\theta\in\mathcal V$ replace the simplex in \eqref{eq:block-map} by $P_\theta-w$. Its dimension is two regardless of the number of active labels. The same Jacobian gives the vertex contribution
\[
\frac{A_\theta}{3}\frac{m(\theta)^2}{r_\theta}
\]
 to $\Vol_3(K)$. Pair contributions are unchanged. The Gaussian pair-flux proof also remains valid, since multiple ties are finitely many rays and have zero area in pair interfaces.

It remains to sum endpoint fluxes of pair arcs. An edge $F$ of $P_\theta$ joins exactly two active signed vectors. Let $e_F$ be its inward unit normal within the polygon plane, so that the other polygon vertices have strictly larger $e_F$-coordinate. Moving from $\theta$ toward $+e_F$ makes those two vertices the unique minimizing pair; all inactive strict inequalities persist. Conversely, any pair arc ending at $\theta$ must arise from such a supporting edge, because its perturbation makes that pair the minimizers of a linear functional on $P_\theta$. Thus the incident pair arcs correspond exactly to polygon edges, with outward conormal $-e_F$.
Let $z_F=a-w$ for a point $a\in F$. Projection onto $W_F$ as in \eqref{eq:block-map} gives
\[
\partial_{n_F}\mu_F=\ip{z_F}{-e_F},\qquad
\sum_{F\text{ edge}}|F|\,\partial_{n_F}\mu_F
=\int_{\partial(P_\theta-w)}\ip{z}{n_{\rm out}}\dd s=2A_\theta.
\]
The last equality is the planar divergence theorem for the field $z\mapsto z$ and does not require the circle center to be inside the polygon. Apply the $d=3,k=2$ divergence identity on all pair arcs, with the zero-endpoint truncation already proved. The total incoming vertex term is $(2A_\theta/3)f_{3,2}(m(\theta))$. Subtracting its volume contribution leaves $(A_\theta/3)q_{3,3}(m(\theta))$, proving the first sum. The map $\theta\mapsto-\theta$ is free on the sphere and sends $P_\theta$ to $-P_\theta$ with unchanged area and radius, giving the second sum and the equality criterion.
\end{proof}

\begin{remark}[A concrete multiple-tie check]\label{rem:tetrahedron}
Take the four tetrahedral normals
\[
\nu_1=(1,1,1)/\sqrt3,\quad\nu_2=(1,-1,-1)/\sqrt3,\quad
\nu_3=(-1,1,-1)/\sqrt3,\quad\nu_4=(-1,-1,1)/\sqrt3.
\]
At each of the six coordinate-axis directions there are four positive active signed normals, forming a square of area $4/3$, with $m=1/\sqrt3$. There are also eight body-diagonal directions with an active equilateral triangle of area $2/\sqrt3$ and $m=1/3$; these are all the directions in $\mathcal V$. To check exhaustion, signed normals are cube vertices. Three of them either share a coordinate face, whose fourth vertex is also active, or, after coordinate reflections, consist of three alternating corners with pairwise Hamming distances two; these give the body-diagonal triangles. A triple containing an opposite pair has common positive projection impossible. Consequently
\begin{equation}\label{eq:tetra-R}
R_{4,3}=\frac83q_{3,3}(1/\sqrt3)+\frac{16}{3\sqrt3}q_{3,3}(1/3).
\end{equation}
The four triangles obtained by selecting three vertices of one square have total area twice the square's area. Replacing strict inactive inequalities by weak ones and summing all those triples would double the square contribution. The polygon formula prevents this error.
\end{remark}

\section{Cosine normals and the conditional comparison}\label{sec:cosine}
For $1\le j\le n$, choose
\[
\nu_j=(\cos((j-1)\pi/n),\sin((j-1)\pi/n),0,\ldots,0).
\]
Their Gram matrix is $\Sigma^{\rm cos}_{ij}=\cos(\pi(i-j)/n)$. Its rank is two for $n\ge2$ and one for $n=1$.
Writing a planar standard Gaussian in polar coordinates shows that its radius $R_2$ and angle are independent, $\E R_2^2=2$, and the distance $\Theta$ of the angle to the nearest perpendicular normal line is uniform on $[0,\pi/(2n)]$. Thus $M=R_2\sin\Theta$ in law, also as recorded in \cite[Proposition 2.1]{Kunisky}. Direct integration gives
\begin{equation}\label{eq:cosine-moment}
\E M^2=\frac{4n}{\pi}\int_0^{\pi/(2n)}\sin^2\theta\dd\theta
=1-\frac n\pi\sin\frac\pi n=c_n.
\end{equation}
The rank-at-most-two identity then yields
\begin{equation}\label{eq:cosine-volume}
\Vol_d(K_{\rm cos})=\kappa_d\frac n\pi\sin\frac\pi n.
\end{equation}
This proves the attainment part of \cref{cor:kunisky} without its conjectural hypothesis.

For normalized spherical measure $\bar\sigma_{d-1}$ define
$Z(\nu,\alpha)=\{u:|\ip{u}{\nu}|\le\sin\alpha\}$, $0\le\alpha\le\pi/2$.
The pointwise identities
\begin{equation}\label{eq:zone-support-superlevel}
\bigcup_i Z(\nu_i,\alpha)=\{u:m(u)\le\sin\alpha\}
=\{u:h_K(u)\ge\cos\alpha\}
\end{equation}
include both endpoints. Thus the union-of-zones profile is exactly the distribution function of the support function $h_K$ on the sphere. In this sense Kunisky's volumetric zone conjecture can be viewed directly as a pointwise extremal statement for the superlevel measures of $h_K$. The Gaussian minimum is a convenient probabilistic encoding of the same spherical data, but it is not needed for the identification \eqref{eq:zone-support-superlevel}.

The nonnegative tail integral and \eqref{eq:radial-moment} give
\begin{equation}\label{eq:zone-integral}
\E M^2=2d\int_0^1s\left[1-\bar\sigma_{d-1}\Bigl(\bigcup_iZ(\nu_i,\arcsin s)\Bigr)\right]\dd s.
\end{equation}
Hence the conjectured pointwise upper bound on zone-union measure by the cosine configuration implies the required lower bound on the Gaussian second moment. The older Fejes T\'oth zone conjecture, proved by Jiang and Polyanskii \cite{JiangPolyanskii}, concerns the threshold at which zones cover the entire sphere; it does not by itself provide the full measure comparison in \eqref{eq:zone-integral}.

For completeness, the stronger stochastic implication of Kunisky's pointwise comparison in a fixed dimension follows by conditioning on the independent Gaussian radius: for $t\ge0$,
\[
\Prob(M\le t)=\E_\rho\left[\bar\sigma_{d-1}\Bigl(\bigcup_i
Z(\nu_i,\alpha(t,\rho))\Bigr)\right],\qquad
\alpha(t,\rho)=
\begin{cases}\pi/2,&0\le\rho\le t,\\ \arcsin(t/\rho),&\rho>t.\end{cases}
\]
For $\rho=0=t$ the chosen value is harmless, since the Gaussian radius has no atom at zero. Moreover $M$ has no atoms: the event $M=t$ is contained in the finite union of events $|\ip{G}{\nu_i}|=t$, each null. Thus comparisons of these distribution functions give the corresponding comparisons of upper tails, even for singular Gram matrices. Integrating tails gives moment comparisons. The chain of sufficient conditions is
\[
\begin{array}{l}
\text{zone-union comparison in dimension }d\\
\qquad\Longrightarrow\ \text{Gaussian stochastic comparison for rank }\le d\\
\qquad\Longrightarrow\ \E M^2\ge c_n\\
\qquad\Longrightarrow\ \eqref{eq:sharp-volume-bound}.
\end{array}
\]
The implications are proved conditionally; none of these conjectural premises is established by this manuscript, and no reverse implication is claimed. The ordinary volume bound itself is exactly the corrected inequality \eqref{eq:corrected-gaussian}.

\section{Concluding remarks}\label{sec:conclusion}
For convex hulls of finitely many great subspheres in arbitrary ambient dimension, the ordinary convex-hull volume is controlled by the same Gaussian minimum that appears in Kunisky's conjecture. The key identity is
\[
\Vol_d(K)=\kappa_d(1-\E M^2)-R_{n,d},\qquad R_{n,d}\ge0.
\]
Under signed affine general position the correction is the explicit nonnegative sum \eqref{eq:R-explicit}, while continuity defines the same total correction for every configuration.

The second intrinsic volume is simpler:
\[
V_2(K)=\pi(d-1-\E M^2).
\]
Thus the Gaussian minimum is encoded exactly in a classical intrinsic volume; the correction $R_{n,d}$ measures the additional loss in passing from this $V_2$ identity to ordinary $d$-dimensional volume for this special class of convex bodies. The proof is dimension-recursive: pairwise Gaussian interfaces feed the $2$-active term, Stokes' theorem transfers the remaining flux to higher active strata, and the scalar functions $q_{d,k}$ record the nonnegative remainder left at each level.

Two additional features isolate the role of degeneracy. The normalized volume and Gaussian law reduce exactly to the span of the normals by \cref{prop:rank-reduction}, yielding equality when that span has dimension at most two. In dimension three, \cref{prop:polygon} gives an explicit formula for arbitrary multiple ties; under general position this reduces to the antipodal signed-triple contribution $L(1-r)^2/(3r)$.

Combining \cref{thm:main} with the $p=2$ portion of Kunisky's conjecture would give the sharp value
\[
\max_K\Vol_d(K)=\kappa_d\frac n\pi\sin\frac\pi n
\]
within the class under consideration, with the cosine configuration as an extremizer. The present argument does not prove the Gaussian-minimum conjecture, does not prove Kunisky's volumetric zone conjecture, and does not establish uniqueness of volume maximizers. The cylindrical/sine-polar description is prior work; the claims made here concern the volume--Gaussian identity and its correction, the associated rank and three-dimensional formulas, and the resulting conditional implication for convex hulls of great subspheres.

\section*{AI assistance and attribution statement}
Generative AI, principally OpenAI's ChatGPT, was used to carry out a substantial majority of the mathematical exploration, conjecture formation, proof development, proof revision, and reproducible diagnostic checking underlying this manuscript. cyclemath formulated and directed the investigation, selected and challenged intermediate claims, requested repeated mathematical audits and corrections, and prepared the manuscript for public release. The designation ``Edited and released by cyclemath'' is intended to record that role without representing the mathematical content as solely or primarily the unaided work of cyclemath. The mathematical arguments and novelty assessment have not received independent verification by a subject-matter expert, and the manuscript has not undergone peer review. The Gaussian-minimum and zone-union conjectures and a classification of all volume maximizers remain outside the proved results.

\begin{thebibliography}{99}
\bibitem{HuangLiXiYe}
Q.~Huang, A.-J.~Li, D.~Xi, and D.~Ye,
\emph{On the sine polarity and the $L_p$-sine Blaschke--Santal\'o inequality},
J. Funct. Anal. \textbf{283} (2022), no.~6, 109571.
\href{https://doi.org/10.1016/j.jfa.2022.109571}{doi:10.1016/j.jfa.2022.109571}.

\bibitem{GordonMinima}
Y.~Gordon, A.~E.~Litvak, C.~Sch\"utt, and E.~Werner,
\emph{Minima of sequences of Gaussian random variables},
C. R. Math. Acad. Sci. Paris \textbf{340} (2005), no.~6, 445--448.
\href{https://doi.org/10.1016/j.crma.2005.02.003}{doi:10.1016/j.crma.2005.02.003}.

\bibitem{GordonMinimum}
Y.~Gordon, A.~E.~Litvak, C.~Sch\"utt, and E.~Werner,
\emph{On the minimum of several random variables},
Proc. Amer. Math. Soc. \textbf{134} (2006), no.~12, 3665--3675.
\href{https://doi.org/10.1090/S0002-9939-06-08453-X}{doi:10.1090/S0002-9939-06-08453-X}.

\bibitem{HenkHernandez}
M.~Henk and M.~A.~Hern\'andez Cifre,
\emph{Intrinsic volumes and successive radii},
J. Math. Anal. Appl. \textbf{343} (2008), no.~2, 733--742.
\href{https://doi.org/10.1016/j.jmaa.2008.01.091}{doi:10.1016/j.jmaa.2008.01.091}.

\bibitem{LitvakSurvey}
A.~E.~Litvak,
\emph{Around the simplex mean width conjecture},
in \emph{Analytic Aspects of Convexity}, Springer INdAM Ser., Springer, 2018, pp.~73--84.
\href{https://doi.org/10.1007/978-3-319-71834-7_5}{doi:10.1007/978-3-319-71834-7\_5}.

\bibitem{MooreCylinders}
M.~Moore,
\emph{Symmetrical intersections of right circular cylinders},
Math. Gaz. \textbf{58} (1974), no.~405, 181--185.
\href{https://doi.org/10.2307/3615957}{doi:10.2307/3615957}.

\bibitem{JiangPolyanskii}
Z.~Jiang and A.~Polyanskii,
\emph{Proof of L\'aszl\'o Fejes T\'oth's zone conjecture},
Geom. Funct. Anal. \textbf{27} (2017), no.~6, 1367--1377.

\bibitem{Kunisky}
D.~Kunisky,
\emph{A revision of Litvak's conjecture on Gaussian minima and a volumetric zone conjecture}, arXiv:2605.02023v1 (2026).
\href{https://arxiv.org/abs/2605.02023}{arXiv:2605.02023}.

\bibitem{Rockafellar}
R.~T.~Rockafellar,
\emph{Convex Analysis}, Princeton Mathematical Series, vol.~28, Princeton University Press, 1970.

\bibitem{Schneider}
R.~Schneider,
\emph{Convex Bodies: The Brunn--Minkowski Theory}, 2nd expanded ed., Encyclopedia of Mathematics and its Applications, vol.~151, Cambridge University Press, 2014.
\end{thebibliography}
\end{document}
